% Квадратичные аппроксимации на последовательных итерациях метода Ньютона:
% вершина очередной параболы становится следующей точкой.
\begin{figure}[htbp]
\centering
\begin{tikzpicture}[scale=1.0, >={Stealth[length=2mm]}]

\def\sx{1.0}
\def\sy{1.0}

% оси
\draw[->] (-0.2,0) -- (6.4,0) node[right, font=\footnotesize] {$x$};
\draw[->] (0,-0.2) -- (0,3.9) node[above, font=\footnotesize] {$f(x)$};

% сама функция: несимметричная "чаша"
\draw[blue!55!black, very thick, domain=0.35:6.0, samples=120, smooth]
  plot (\x, {0.42*(\x-3.4)*(\x-3.4) + 0.35 + 0.9/(\x+0.3)});
\node[blue!55!black, font=\footnotesize, anchor=west] at (5.4,3.0) {$f$};

% три квадратичные аппроксимации, вершины в x_0 > x_1 > x_2 -> x_*
% аппроксимация в x_0 = 1.0 (вершина следующей точки x_1 = 2.0)
\draw[orange!75!black, thick, domain=0.75:3.6, samples=60, smooth]
  plot (\x, {1.15*(\x-2.0)*(\x-2.0) + 1.05});
% аппроксимация в x_1 = 2.0 (вершина x_2 = 2.95)
\draw[orange!60!black, thick, domain=1.9:4.2, samples=60, smooth]
  plot (\x, {0.80*(\x-2.95)*(\x-2.95) + 0.62});
% аппроксимация в x_2 = 2.95 (вершина близко к x_*)
\draw[orange!45!black, thick, domain=2.3:4.6, samples=60, smooth]
  plot (\x, {0.55*(\x-3.35)*(\x-3.35) + 0.48});

% отметки точек на оси
\foreach \p/\lab in {1.0/{x_k}, 2.0/{x_{k+1}}, 2.95/{x_{k+2}}} {
  \draw[gray, dashed] (\p,0) -- (\p,{0.42*(\p-3.4)*(\p-3.4) + 0.35 + 0.9/(\p+0.3)});
  \fill[orange!75!black] (\p,0) circle (1.5pt);
}
\node[font=\scriptsize, anchor=north] at (1.0,-0.08) {$x_k$};
\node[font=\scriptsize, anchor=north] at (2.05,-0.08) {$x_{k+1}$};
\node[font=\scriptsize, anchor=north] at (3.05,-0.08) {$x_{k+2}$};

\end{tikzpicture}
\caption{Метод Ньютона как последовательность квадратичных аппроксимаций. В
точке $x_k$ строится парабола $\psi_k$, совпадающая с $f$ по значению,
градиенту и гессиану; её вершина даёт $x_{k+1}$, то есть единичный шаг
$\alpha_k = 1$ вдоль ньютоновского направления. Чем ближе к минимуму, тем
теснее парабола прилегает к $f$ — отсюда квадратичная скорость.}
\label{fig:newton-quadratic}
\end{figure}
